\problem{One-Round Random-Priority Conflict Filter}

A system has a set of tasks. Some pairs of tasks conflict and cannot be selected simultaneously. This conflict relation is represented by a simple undirected graph \(G=(V,E)\): each vertex is a task, and each edge means that the two corresponding tasks conflict.

The system performs the following randomized filtering step:
\begin{enumerate}
    \item Each vertex \(v\) independently chooses a priority \(p(v)\) uniformly at random from \([0,1]\).
    \item Vertex \(v\) is selected if its priority is smaller than the priority of every neighbor of \(v\). Here a smaller numerical priority means a higher priority.
\end{enumerate}

Because the priorities are drawn from a continuous distribution, all priorities are distinct with probability \(1\). Let \(S\) be the set of selected vertices.

Answer the following questions:
\begin{enumerate}[(a)]
    \item Prove that \(S\) is always an independent set.
    \item For a fixed vertex \(v\), let \(\deg(v)\) be its degree. What is the probability that \(v\) is selected?
    \item Let \(X=|S|\). Use indicator random variables to compute \(\mathbb{E}[X]\), and simplify the answer when \(G\) is \(d\)-regular.
\end{enumerate}

\solution{
\textbf{(a)}
Suppose for contradiction that \(S\) is not an independent set. Then there is
an edge \(\{u,v\}\in E\) such that both \(u\) and \(v\) are selected.

Since \(u\) is selected and \(v\) is its neighbor, we must have
\[
    p(u)<p(v).
\]
Since \(v\) is selected and \(u\) is its neighbor, we must also have
\[
    p(v)<p(u).
\]
This is impossible. Hence no two adjacent vertices can both be selected, and
\(S\) is always an independent set.

\textbf{(b)}
Vertex \(v\) is selected iff its priority is the smallest among the
\(\deg(v)+1\) vertices consisting of \(v\) and all its neighbors. These
priorities are independent and drawn from the same continuous distribution, so
each of the \(\deg(v)+1\) vertices is equally likely to have the smallest
priority. Therefore
\[
    \Pr[v\in S]=\frac{1}{\deg(v)+1}.
\]

\textbf{(c)}
For each vertex \(v\in V\), define the indicator random variable
\[
    X_v =
    \begin{cases}
    1, & \text{if } v\in S,\\
    0, & \text{otherwise}.
    \end{cases}
\]
Then
\[
    X=|S|=\sum_{v\in V}X_v.
\]
By linearity of expectation and part (b),
\[
    \mathbb{E}[X]
    =
    \sum_{v\in V}\mathbb{E}[X_v]
    =
    \sum_{v\in V}\Pr[v\in S]
    =
    \sum_{v\in V}\frac{1}{\deg(v)+1}.
\]

If \(G\) is \(d\)-regular, then \(\deg(v)=d\) for every \(v\), and therefore
\[
    \mathbb{E}[X]
    =
    \frac{|V|}{d+1}.
\]
}
\newpage
